\documentclass{assignment}

\course{ECO 120-04}
\name{Lucas Reddinger}
\date{Friday 16 September 2022}
\doctitle{Assignment 2: Production possibility frontier}

\begin{document}
\RaggedRight

\beginassignment{}

You may work with other students on this assignment. However, most of the work here should be your own, so that you may obtain valuable feedback on your personal progress.

\emph{Due Friday 23 September.} Please submit hardcopy at the beginning of class (11:00 a.m.), or if you prefer, under the door of Wimberly Hall 339C by 10:50 a.m.

\section{An economy of kale and dental floss}

Your small country produces two goods, kale and dental floss.

When production is efficient, you can produce the following combinations of kale and dental floss:

\begin{tabular}{l>{\raggedleft\arraybackslash}p{0.9in}>{\raggedleft\arraybackslash}p{1.3in}}
\toprule
Point & Kale (tons) & Dental floss (km) \\
\midrule
$A$ &   0 & 1,000 \\
$B$ & 200 &   900 \\
$C$ & 400 &   700 \\
$D$ & 600 &   400 \\
$E$ & 800 &     0 \\
\bottomrule
\end{tabular}

These points are connected with straight line segments; altogether these comprise the production possibility frontier.

\clearpage

\begin{enumerate}

\item Please graph the production possibility frontier (PPF), with kale on the horizontal axis and dental floss on the vertical axis. Use a scale of 0 to 1,200 on the horizontal axis and 0 to 1,000 on the vertical axis. Properly label the graph by labelling both the axes and the curve.

\vspace{12pt}
\begin{center}
\begin{tikzpicture}[x=4mm, y=4mm]
\clip(3,3) rectangle (42,27);
\draw[step=4mm, very thin, black!20!white] (0,0) grid (60,120);
\draw[step=20mm, thick, black!40!white] (0,0) grid (60,120);
\end{tikzpicture}
\end{center}
\vspace{24pt}

\item The PPF you drew above should consist of four line segments---$\overline{AB}$, $\overline{BC}$, $\overline{CD}$, and $\overline{DE}$. For each segment, please calculate the cost of one ton of kale. Ensure that you specify the units of the cost. Please show your work. If you use a formula, please write it clearly upon its first use.

\clearpage

\item Now for each segment, please calculate the cost of one kilometer of dental floss. Again, please show your work and ensure that you specify the units of the cost.

\vfill

\item What is the cost of the 500th ton of kale? Please specify units.

\vspace{5\baselineskip}

\item What is the cost of the 921st km of dental floss? Please specify units.

\vspace{5\baselineskip}

\end{enumerate}

\clearpage

\section{Technological progress}

New irrigation technology allows you to produce more kale. Now efficient production is given by the following table:

\begin{tabular}{l>{\raggedleft\arraybackslash}p{0.9in}>{\raggedleft\arraybackslash}p{1.3in}}
\toprule
Point & Kale (tons) & Dental floss (km) \\
\midrule
$A'$ &     0 & 1,000 \\
$B'$ &   300 &   900 \\
$C'$ &   600 &   700 \\
$D'$ &   900 &   400 \\
$E'$ & 1,200 &     0 \\
\bottomrule
\end{tabular}

Once again, these points are connected with straight line segments, comprising the new PPF.

\begin{enumerate}

\item Please draw this new PPF on your original graph, labeled as PPF'.

\item With your new irrigation technology, what is the cost of the 500th ton of kale? Please show your work and specify units.

\vfill

Did the cost increase, decrease, or remain the same compared to the cost with the old technology?

\vspace{4\baselineskip}

\item With your new irrigation technology, what is the cost of the 921st km of dental floss? Please show your work and specify units.

\vfill

Did the cost increase, decrease, or remain the same compared to the cost with the old technology?

\vspace{4\baselineskip}

\end{enumerate}

\end{document}
